% Solutions to the exercises of Chapter 18.
\begin{solutions}
  \item With $H = [\,I_3\ \ 0\ \ {-\hat R\sx{\ell_b}}\ \ 0\ \ 0\,]$ and the block-diagonal $P$, only the position and attitude blocks contribute:
        \[
          S = \sigma_p^2I + \sigma_\theta^2\hat R\sx{\ell_b}\sx{\ell_b}\T\hat R\T + \sigma^2I = \alpha\,I - \sigma_\theta^2\,(\hat R\ell_b)(\hat R\ell_b)\T,
          \qquad \alpha = \sigma_p^2 + \sigma^2 + \sigma_\theta^2\|\ell_b\|^2,
        \]
        using $\sx{\ell}\sx{\ell}\T = -\sx{\ell}^2 = \|\ell\|^2I - \ell\ell\T$. The gain blocks are $K_p = \sigma_p^2S\inv$, $K_\theta = \sigma_\theta^2(-\hat R\sx{\ell_b})\T S\inv = \sigma_\theta^2\sx{\ell_b}\hat R\T S\inv$,
        and zero for velocity and biases. Since $S\inv(\hat R\ell_b) \propto \hat R\ell_b$ and $\sx{\ell_b}\hat R\T\hat R\ell_b = 0$, the attitude gain simplifies to
        \[
          K_\theta = \frac{\sigma_\theta^2}{\sigma_p^2 + \sigma^2 + \sigma_\theta^2\|\ell_b\|^2}\,\sx{\ell_b}\hat R\T .
        \]
        Its range is the range of $\sx{\ell_b}$: the update corrects only the attitude components perpendicular to the lever arm, never the rotation
        about $\ell_b$ itself, which does not move the antenna. Velocity and biases are not corrected here because this prior has no cross-covariances;
        after some propagation they would be.

  \item With $h(x) = \|p - b_w\|$ and $\hat\rho = \|\hat p - b_w\|$, perturbing $p = \hat p + \delta p$ gives $h \approx \hat\rho + u\T\delta p$ with $u = (\hat p - b_w)/\hat\rho$, so
        \[
          H = [\,u\T\ \ 0\ \ 0\ \ 0\ \ 0\,] \in \R^{1\times15} .
        \]
        Its nullspace is 14-dimensional: the two position directions perpendicular to the line of sight (moving on the sphere around the beacon), and
        all velocity, attitude and bias errors. One range observes one number; position perpendicular to the line of sight becomes observable only
        through several beacons or through motion and the correlations it builds.

  \item Simulate a constant-velocity target (acceleration density $\sqrt{0.5}$, position measurements at 10~Hz with unit variance) and run the
        filter with $Q$ divided by 100 and a $\chi^2_{0.95}(1)$ gate. Typical results, per 30-second window:
        \begin{center}\small
        \begin{tabular}{@{}l l l@{}}
        \toprule
        filter & rejection rate & RMS position error \\
        \midrule
        correct $Q$, gated & 2--8\% & 0.4--0.65~m \\
        $Q/100$, gated & 79\%, then 100\% & 26, 126, 243, \dots, 3246~m \\
        $Q/100$, not gated & 0\% & 0.8--1.7~m \\
        \bottomrule
        \end{tabular}
        \end{center}
        The overconfident filter predicts innovations that are far too small, so good measurements fail the gate; each rejection removes the
        correction the filter needed, the error grows, more measurements are rejected, and within half a minute the filter ignores all data and
        dead-reckons. Its reported $\sigma$ (about 200~m at the end) stays far below its actual error. Without the gate the same mistuned filter is
        merely suboptimal: its error stays bounded. A rejection rate far above the design value is the alarm signal (\cref{sec:gating}).

  \item At rest, the accelerometer measures $\hat R\T(-g_w) + b_a$, and to first order a tilt error and a horizontal bias change this reading in the same way:
        in the velocity error dynamics both enter through the horizontal specific-force error, so from accelerations alone only their combination is
        observable (\cref{ch:strapdown}). With a pose sensor, the position measurements make the velocity error and its rate observable (two
        differentiations, as in \cref{sec:observability-linear}), while the attitude measurement observes $\delta\theta$ directly. The rate of the velocity
        error is (tilt term) $+$ (bias term); with the tilt term known from the attitude measurement, the bias term, and hence the horizontal bias, is
        determined. The vertical bias was observable already (it changes the measured magnitude of gravity, which the position measurements expose),
        and the gyro bias is observable from the drift of the attitude between measurements. With a pose sensor, every IMU state is observable,
        which is why the covariances converge in Milestone 6.

  \item For the optimal gain, $P_{\mathrm{prior}} - P_{\mathrm{post}} = KHP = KSK\T$ (\cref{thm:kf-update}), and $KSK\T \succeq 0$ because it is a congruence of $S \succ 0$
        (\cref{prop:psd}). In floating point the simple update $P_{\mathrm{post}} = (I - KH)P$ is neither symmetric nor guaranteed PSD (Exercise 12.3), and
        $P_{\mathrm{prior}} - P_{\mathrm{post}}$ then has rounding-level asymmetry and can have slightly negative eigenvalues, notably in directions the
        measurement does not observe, where the true difference is exactly zero and the computed one is $\pm u\|P\|$. A test must therefore symmetrize the
        difference and allow a tolerance: $\lambda_{\min}\big(\tfrac12(D + D\T)\big) \ge -c\,u\,\|P\|$ with a small constant $c$. Using the Joseph form and
        symmetrizing after each update makes $P_{\mathrm{post}}$ itself PSD, but the difference still needs the tolerance.
\end{solutions}
